% Nouveau EGA V (cote 25): archive cote 25, site page 140 (PDF page 141).
% English translation of the checked French annotated master.
% AI-translated by GPT-6 Astra (OpenAI, ChatGPT Pro mode).
\scanpage{25/140}
\var{\margin{\uline{new EGA} \uuline{V}}}{\placedmargin{\uline{new EGA} \uuline{V}}}{marginal note placed}{ega25_25_140-v0001}

IV 16

A new n$^{\text{o}}$: \uline{Singular and supersingular set of a function},
\uline{and differential criteria}.

This n$^{\text{o}}$ will be used in section 20 on hyperplane sections, but its
natural place seems to me to be in section 16.

\uline{Definition} 1 Let X be a regular prescheme, $\hand{\phi}$ a section of $\underline{O}_X$. A
point $x \hand{\in} X$ is called a \uline{singular} \var{\xout{\illegible{3}}}{}{typewriter strike-out}{ega25_25_140-v0002} \uline{zero} for $\hand{\phi}$
if we have $\hand{\phi}_x \hand{\in} \underline{m}_x^2$, a \uline{supersingular} zero if it is a singular zero
and if in addition the element of $\underline{m}_x^2/\underline{m}_x^3 \hand{\simeq} \mathrm{Sym}(\underline{m}_x/\underline{m}_x^2)$ that it defines,
interpreted as a quadratic form on the dual $\underline{t}_x$ of $\underline{m}_x/\underline{m}_x^2$ over k(x),
is a \var{\deleted{non-}}{}{deletion}{ega25_25_140-v0003}degenerate form. (A singular zero that is not supersingular
is sometimes called an \uline{ordinary singular zero}).

\uline{Remarks} 2 If $x \hand{\in} V(\hand{\phi})$\hand{,} then x is a \var{\inserted[typed]{non-}}{non-}{insertion}{ega25_25_140-v0004}singular zero of $\hand{\phi}$ iff
$\hand{\phi}_x \hand{\neq} 0$ and x is a nonsingular, i.e. regular, point of $V(\hand{\phi})$,
i.e. iff x is a regular point of $V(\hand{\phi})$ and $V(\hand{\phi}) \hand{\neq} X$ in a neighborhood of
x. \var{\deleted{One can prove that if x is singular for $\hand{\phi}$, it is ordinary
singular iff $\hand{\phi}_x \hand{\neq} 0$ and the multiplicity of $\underline{O}_{V(\hand{\phi}),x}$ is equal to 2
(or, as one also sometimes says, x is an \uline{ordinary singular point}
\uline{of the prescheme} $V(\hand{\phi})$); thus an $x \hand{\in} V(\hand{\phi})$ is not supersingular \var{\inserted[typed]{for $\hand{\phi}$}}{for $\hand{\phi}$}{insertion}{ega25_25_140-v0007} iff
$\hand{\phi}_x \hand{\neq} 0$ and the multiplicity of $O_{V(\hand{\phi}),x}$ is $\hand{\leq} 2$.}}{}{deletion}{ega25_25_140-v0006}

\uline{Definition} 3. Let X be a prescheme smooth over a field k, $\hand{\phi}$ a
section of $\underline{O}_X$, $x \hand{\in} V(\hand{\phi})$. We say that x is a \uline{geometrically}
\uline{singular} (resp. \uline{geometrically \unsure{supersingular}}) zero of $\hand{\phi}$ relative to k,
if for every extension k' of k and every point \hand{$\xi$} of X with values in
k\hand{,} localized at x, the \var{\xout{\illegible{1}}}{}{typewriter strike-out}{ega25_25_140-v0009} corresponding point x' of $X_{k'}$ is a
\var{\xout{\unsure{geom}}}{}{typewriter strike-out}{ega25_25_140-v0010} singular (resp. supersingular) zero of $\hand{\phi}_{k'}$.

\mspage{\hand{140}}
\tnote{25-140-1}{The source abbreviation n with raised o means a number or subsection. The canceled doubtful fragments régxxx and dexx have no separate English counterpart; their readings remain recorded here. The canceled geom abbreviates geometrically. The definition retains k, without a prime, for the field in which the point has values, as in the checked French text.}
